% !TEX root = ../paper.tex

\section{Results}
\label{sec:empirical_results}

\subsection{Permit Activity After Redistricting}
\label{sec:empirical_results_permits}

I first ask whether permit activity changes when the same location is assigned to a different alderman.
Chicago's 2015 ward redistricting moved some census blocks toward wards represented by more stringent aldermen and others toward wards represented by more lenient aldermen.
Nearby blocks that remained in their original ward provide the comparison group.
The design therefore holds the physical location fixed while changing its assigned regulatory environment.
The primary analysis sample keeps boundaries for which both the origin- and destination-ward aldermen serving in 2014 remained in office when the new map took effect,
since newly elected aldermen's development preferences could be endogenous to prior development activity in the ward.
I also require each block to have at least one high-discretion permit during the five years before redistricting.
Without that restriction, any later permit on a block with no previous activity would mechanically appear as an increase.

The outcome is the number of high-discretion permits recorded on each census block in each year from 2010 through 2020.
The outcome therefore measures when projects observed in the permit data first appeared in the process.
These are the high-discretion permits defined in Section~\ref{subsec:permit_pipeline}.

I estimate a Poisson pseudo-maximum likelihood event study that treats moves toward greater stringency and greater leniency as opposite changes:
%
\begin{equation}
\label{eq:event_study_permits}
\begin{aligned}
\log \mathbb{E}[Y_{by}\mid\cdot] ={}&
\mu_b + \lambda_{g(b),y}
+ \sum_{t \neq -1}
\beta_t Z_b \mathbf{1}\{y-2015=t\}.
\end{aligned}
\end{equation}
%
Here $Y_{by}$ is the permit count for block $b$ in calendar year $y$, $\mu_b$ is a block fixed effect, and $\lambda_{g(b),y}$ is a ward-pair-by-year fixed effect for the ward pair $g(b)$ containing block $b$.
Event time is $t=y-2015$, with $t=-1$ omitted.
Let $\Delta S_b$ be the destination-minus-origin difference in standardized aldermanic stringency for block $b$.
For blocks that change wards, $Z_b=\operatorname{sign}(\Delta S_b)$; it equals zero for blocks that remain in the same ward.
I classify moves using the 2014 incumbents in each block's origin and destination wards, with scores estimated only from permits filed between 2006 and 2014.
Using the 2014 incumbents separates ward reassignment from changes in representation associated with the 2015 elections.
This specification imposes equal-and-opposite effects in logs for assignments toward more stringent and more lenient aldermen.
Panel A of Figure~\ref{fig:permit_event_study} plots the coefficients $\beta_t$.
The pre-trends are not distinguishable from zero.
The pooled 2015--2020 estimate is $-0.126$ (SE $=0.049$), equivalent to an 11.8\% reduction in the annual number of high-discretion permits associated with a move toward greater stringency.
Panels B and C relax the symmetry restriction by estimating the two directions separately.
The decline in permitting activity comes mostly from moves toward more stringent aldermen, which reduce high-discretion permits by 20.6\%; blocks assigned toward more lenient aldermen experience a 6.9\% increase.
\onlineappendixlink{Table C.2}{tab:permit_event_study_did} reports the corresponding pooled estimates without the symmetry restriction and shows the results are nearly identical.

This asymmetry suggests that aldermanic discretion operates through uniquely stringent aldermen blocking activity that would otherwise occur, rather than uniquely lenient aldermen allowing activity that formal city regulations would otherwise block.
This distinction matters for the likely aggregate effect of informal regulatory discretion on housing supply because it is possible ex ante that a few particularly pro-development aldermen approve more housing than a centralized system would.
The results instead suggest that informal aldermanic discretion reduces housing supply relative to a more centralized counterfactual.

\begin{figure}[htbp]
    \centering
    \includegraphics[width=0.62\textwidth]{../tasks/run_event_study_permit/output/permit_event_study_high_discretion_stable_signed_500ft.pdf}

    \vspace{0.3em}
    \includegraphics[width=0.96\textwidth]{../tasks/run_event_study_permit/output/permit_event_study_high_discretion_stable_separate_500ft.pdf}
    \caption{Permit Activity After the 2015 Ward Redistricting}
    \figurenotes{The figure compares annual high-discretion permit activity on blocks reassigned to more- or less-stringent aldermen in 2015 with nearby blocks that stayed in their wards.
    The 2010--2020 sample includes blocks within 500~ft of ward boundaries; \onlineappendixlink{A.3}{subsec:event_study_panel} details sample construction.
    Panel A combines both reassignment directions using Equation~\eqref{eq:event_study_permits}; Panels B and C allow the effects of the two directions to differ.
    Points show effects in log points relative to 2014; subtitles report pooled 2015--2020 coefficients on the same scale.
    Panel A's pooled coefficient of $-0.126$ corresponds to an 11.8\% reduction in annual permits for assignment toward greater stringency.
    Shaded areas are 95\% confidence intervals based on standard errors clustered by ward pair.
    The test that the pre-2015 effects in Panel A are jointly zero gives $p=0.395$.
    * $p\leq0.10$, ** $p\leq0.05$, *** $p\leq0.01$.}
    \label{fig:permit_event_study}
\end{figure}

\subsection{Density at Ward Boundaries}
\label{sec:empirical_results_density}

I next examine the individual new construction projects described in Section~\ref{subsec:new_construction_data} on opposite sides of the same ward boundary, moving the unit of analysis from census blocks to projects.
Projects within 500~ft of the same boundary segment face similar neighborhood amenities and demand conditions, but different aldermen.
Following \citet{bayer_unified_2007} and \citet{kulka_under_2026}, I divide signed distance to the boundary into 100~ft bands and compare projects in each band with projects in the nearest band on the less-stringent side.
I estimate the following specification:
%
\begin{equation} \label{eq:segment_fe}
\begin{aligned}
    \log(Y_{isy}) ={}&
    \sum_{k \in \mathcal{K},\,k \neq [-100,0)}
    \kappa_k \mathbf{1}\{d_{is} \in k\}
    + X_i'\theta \\
    &\quad
    + \zeta_{z(i,y)}
    + \delta_s
    + \gamma_y
    + \varepsilon_{isy}.
\end{aligned}
\end{equation}
%
Here $Y_{isy}$ is FAR or DUPAC for project $i$, assigned to boundary segment $s$ and built in calendar year $y$.
The running variable $d_{is}$ is negative on the less-stringent side and positive on the more-stringent side.
The set $\mathcal{K}$ contains ten 100~ft bands from $-500$ to 500~ft.
The band from $-100$ to 0~ft is omitted, so $\kappa_{[0,100)}$ compares projects immediately across the boundary.
The vector $X_i$ contains the White population share, the Black population share, median household income, the bachelor's-degree share, and the homeownership rate on the project's side of the boundary.
The model also includes the zoning group in force when the project was built, boundary-segment fixed effects, and construction-year fixed effects.
Standard errors are clustered by ward pair.

Figure~\ref{fig:rd_fe_plots} plots the distance-band estimates.
The omitted band is normalized to zero, and the subtitle in each panel reports $\kappa_{[0,100)}$.
For all new construction, FAR is about 9\% lower and dwelling units per acre about 14\% lower immediately across the boundary on the more-stringent side.
Among multifamily projects, FAR is about 13\% lower and dwelling units per acre about 17\% lower.
All four differences are statistically distinguishable from zero at the 10\% level or better.
The similar results for FAR and dwelling units per acre indicate that the difference is present in both floor space and the number of homes placed on a parcel.

\begin{figure}[htbp]
    \centering
    \includegraphics[width=0.96\textwidth]{../tasks/density_main_results/output/density_rd.pdf}
    \caption{Development Density at Ward Boundaries}
    \figurenotes{Left column: all residential new construction; right column: new multifamily construction.
    Top row: floor-area ratio; bottom row: dwelling units per acre, both in log differences.
    Samples include new construction within 500~ft of ward boundaries with both FAR and DUPAC available.
    Each point is the coefficient for a 100~ft distance band from Equation~\eqref{eq:segment_fe}.
    The nearest band on the less-stringent side, from $-100$ to 0~ft, is omitted and normalized to zero.
    Subtitles report the coefficient for the nearest band on the more-stringent side, from 0 to 100~ft.
    Shaded areas are 95\% confidence intervals based on standard errors clustered by ward pair.
    Negative distance indicates the less-stringent side and positive distance the more-stringent side.
    * $p<0.10$, ** $p<0.05$, *** $p<0.01$.}
    \label{fig:rd_fe_plots}
\end{figure}

As a falsification test, I repeat the same comparison at artificial cutoffs 1,000~ft inside either side of the true ward boundary.
These locations do not mark a change in aldermanic authority.
\onlineappendixlink{Figure D.1}{fig:rd_placebo_shifts} shows little evidence of corresponding breaks at these artificial cutoffs.
None of the eight placebo estimates is statistically distinguishable from zero at the 10\% level.
The contrast with the true boundary makes it less likely that the main estimates reflect a general spatial pattern within wards.
\onlineappendixlink{D}{sec:appendix_density} presents the placebo figures along with continuity tests for fixed location characteristics, restrictions based on boundary geometry, and donut checks.

\subsection{Evidence on Rents and Sale Prices}
\label{sec:empirical_results_prices}

The permit and density results indicate that more stringent aldermen reduce development activity and the density of completed housing.
The model in Section~\ref{sec:model} indicates that this should feed into housing costs, the most directly policy-relevant outcome given Chicago's current rent and price growth.
For rents and sales, I estimate
%
\begin{equation}
\label{eq:price_rd}
    \log(Y_{isp}) =
    \sum_{k \in \mathcal{K},\,k \neq [-100,0)}
    \kappa_k \mathbf{1}\{d_{is} \in k\}
    + X_i'\theta
    + \lambda_{t(i)}
    + \delta_{sp}
    + \varepsilon_{isp},
\end{equation}
%
where $Y_{isp}$ is the real listed rent or real sale price for observation $i$ near boundary segment $s$ in calendar period $p$.
As in the density analysis, $\mathcal{K}$ contains ten 100~ft distance bands and the band from $-100$ to 0~ft is omitted.
The coefficient $\kappa_{[0,100)}$ therefore compares observations in the two bands immediately across the boundary.
The fixed effects compare observations near the same boundary segment in the same month for rents and the same quarter for sales.
$X_i$ includes unit or property characteristics and distances to nearby schools, parks, major roads, CTA stops, and Lake Michigan.
The term $\lambda_{t(i)}$ denotes property-type fixed effects: building type reported in rental listings for rents and Cook County Assessor property class for sales.
Standard errors are clustered by ward pair.

Figure~\ref{fig:price_boundary_bins} plots these distance-band estimates.

Listed rents are 1.8\% higher on the more-stringent side relative to the less-stringent side, although the estimate is imprecise.
Home sale prices are 3.0\% higher on the more-stringent side relative to the less-stringent side.
I treat the sales result as suggestive because transactions are less frequent near any particular boundary and their composition is harder to hold fixed.
\onlineappendixlink{E}{sec:appendix_price_rd} shows that the result disappears at artificial cutoffs inside either ward and is similar when I restrict the sample to straight boundaries or exclude observations immediately next to the boundary.

\begin{figure}[htbp]
    \centering
    \includegraphics[width=0.98\textwidth]{../tasks/price_boundary_results/output/price_boundary_main.pdf}
    \caption{Listed Rents and Home Sale Prices at Ward Boundaries}
    \figurenotes{The left panel uses monthly asking rents by floorplan from 2014--2022, with segment-by-month fixed effects.
    The right panel uses 58,468 sales of detached houses, townhouses, and apartment buildings with two to six units from 2006--2022, with segment-by-quarter fixed effects.
    Both panels control for housing characteristics and distances to nearby amenities.
    Rental controls include listing building-type fixed effects; sales controls include Cook County Assessor property-class fixed effects.
    Each point is the estimated log difference between a 100~ft distance band and the band immediately inside the less-stringent side, which is normalized to zero.
    Shaded areas are 95\% confidence intervals.
    Subtitles report the comparison between the two bands adjacent to the boundary.
    Standard errors are clustered by ward pair. * $p<0.10$, ** $p<0.05$, *** $p<0.01$.}
    \label{fig:price_boundary_bins}
\end{figure}
